% =====================================================================
% Explicit VOA content supplementing universal_anomaly_four_climax_simultaneously_2026_04_22.tex
% Five load-bearing computations: H_{Delta_5} comparison target, master trace on K3,
% Theorem B chain-level on H_1 and V_1(sl_2), DVV 1997 Sym^N(K3) derivation,
% nine Heegner arithmetic VOAs.
%
% Intended as an \input supplement to the main standalone; carries the
% same preamble assumptions (same LaTeX class, same \cH, \cO, \bfH macros).
%
% Author: Raeez Lorgat. 2026-04-23.
% =====================================================================

% ==================================================
\section{Explicit VOA content: five load-bearing computations}
\label{sec:voa-explicit}
% ==================================================

Each of the comparison vertices of Theorem~\ref{thm:main} carries a
concrete vertex-operator-algebra presentation with explicit OPEs,
explicit $\kappa$-invariants, and explicit Koszul-dual toy models on
the reference curve $C$. The computations below are evidence and input
for the conditional bridges; they do not by themselves prove the
five-vertex comparison conjecture.

% ==================================================
\subsection{$\bfH_{\Delta_5}$ comparison target: $\mathfrak g_{\Delta_5}$ generators and the Borcherds denominator $1/\Delta_5$}
\label{subsec:voa-H-Delta5-canonical}
% ==================================================

The K3 chiral bialgebra $\bfH_{\Delta_5}$ of
Theorem~\ref{thm:climax-III} is the conformal-block lift of the K3
Borcherds--Kac--Moody superalgebra $\mathfrak g_{\Delta_5}$ (Gritsenko--Nikulin 1998) via a lattice vertex
algebra on the Mukai lattice $L = H^\bullet_{\mathrm{even}}(K3,\Zbb)$
of signature $(4,20)$. The Frenkel--Kac lift transports every
BKM generator to a vertex-operator residue.

\paragraph{Lattice vertex algebra on the Mukai lattice.}
The Mukai lattice $L \cong \Zbb^{4,20}$ carries the positive-definite
$c_+(L) = 4$ summand generated by
$\beta^{(1)},\beta^{(2)},\beta^{(3)},\beta^{(4)}$ (class of a fixed
polarisation plus $H^0 \oplus H^4$ together with two transverse
$(1,1)$-classes), and the negative-definite $c_-(L) = 20$ summand
generated by transverse $(1,1)$ and $H^{2,0}+H^{0,2}$ directions. Let
$\phi^{(a)}(z)$, $a = 1,\ldots,24$, be $24$ free bosons of conformal
weight one realising the Mukai pairing:
\begin{equation}
\label{eq:voa-muk-ope}
  \phi^{(a)}(z)\,\phi^{(b)}(w) \;\sim\; g^{ab}_{\mathrm{Muk}}\log(z-w),
  \qquad g^{ab}_{\mathrm{Muk}} \text{ Mukai form of signature }(4,20).
\end{equation}
For $\alpha \in L$, the lattice vertex operator is
$V_\alpha(z) = {:}\exp(\alpha\cdot\phi(z)){:}$ with OPE (Borcherds 1986, \S5)
\begin{equation}
\label{eq:voa-lattice-vo-ope}
  V_\alpha(z)\,V_\beta(w)
    \;=\;
    \epsilon(\alpha,\beta)\,(z-w)^{\langle\alpha,\beta\rangle}
    \,{:}V_\alpha(z)V_\beta(w){:} \;+\; (\text{regular}),
\end{equation}
with $\epsilon$ the Borcherds sign cocycle. The Fock module
$\cF_{K3} = \mathrm{Sym}\bigl(\bigoplus_{n\geq 1}L\otimes t^{-n}\bigr)
\otimes \Cbb[L]$ is the universal Heisenberg module on which every
vertex operator $V_\alpha$ acts.

\paragraph{Real and imaginary simple roots of $\mathfrak g_{\Delta_5}$.}
The BKM superalgebra $\mathfrak g_{\Delta_5}$ has a rank-three
Lorentzian primitive Cartan lattice, not an even unimodular lattice of
signature $(2,1)$; even unimodular indefinite lattices require
signature difference divisible by $8$. In the normalization used here
its real simple roots $\delta_1,\delta_2,\delta_3$ have Gram matrix
\begin{equation}
\label{eq:voa-BKM-Gram}
  G_{\mathrm{BKM}} \;=\; \begin{pmatrix} 2 & -2 & -2 \\ -2 & 2 & -2 \\ -2 & -2 & 2 \end{pmatrix}
  \;=\; 2\,\mathrm{diag}(1,1,1) \;-\; 2(J - I), \qquad \det G_{\mathrm{BKM}} = -32,
\end{equation}
and imaginary simple roots $\alpha^{\mathrm{im}} = (n,\ell,m) \in \Lambda_{\Delta_5}$
(with $4nm - \ell^2 \geq 0$ a Humbert discriminant invariant) of
multiplicity
\begin{equation}
\label{eq:voa-imag-mult}
  \mathrm{mult}(\alpha^{\mathrm{im}}) \;=\; c_{K3}(4nm - \ell^2),
\end{equation}
with $c_{K3}(D)$ the Fourier coefficient of $D$-th order in the K3
elliptic genus generator
$\phi^{K3}_{0,1}(\tau,z) = 2\phi_{0,1} - (E_2/6)\phi_{-2,1}$
(Eichler--Zagier 1985; Dabholkar--Murthy--Zagier 2012).

\paragraph{Real-root Chevalley generators as Heisenberg vertex operators.}
For each real simple root $\delta_i$, the Chevalley generators
$e_i,f_i,h_i$ are realised as bare Heisenberg currents on the
negative-definite sub-lattice directions of the Mukai lattice:
\begin{equation}
\label{eq:voa-real-root-generators}
  e_i(z) \;=\; V_{\delta_i}(z) \;=\; {:}e^{\delta_i\cdot\phi(z)}{:}, \qquad
  f_i(z) \;=\; V_{-\delta_i}(z), \qquad
  h_i(z) \;=\; \delta_i\cdot\partial\phi(z).
\end{equation}
The Chevalley commutator $[e_i,f_j] = \delta_{ij}h_i^\vee$ follows from
the $(z-w)^{\langle\delta_i,-\delta_j\rangle} = (z-w)^{-2}$ double pole
of the lattice-vertex-operator OPE, with the residue extracting the
Cartan generator. The Serre relations at $\delta_i,\delta_j$ with
$\langle\delta_i,\delta_j\rangle = -2$ follow from the two-term Serre
inclusion $\mathrm{ad}(e_i)^{1 - a_{ij}}(e_j) = \mathrm{ad}(e_i)^3(e_j) = 0$
witnessed by the fourth-order pole vanishing in the triple OPE
$V_{\delta_i}V_{\delta_i}V_{\delta_i}V_{\delta_j}$.

\paragraph{Imaginary-root Chevalley generators as screening operators.}
For each imaginary simple root $\alpha^{\mathrm{im}}$, the Chevalley
generator $e^{\mathrm{im}}_\alpha$ is a Feigin--Frenkel screening
operator
\begin{equation}
\label{eq:voa-imag-root-screening}
  S_{\alpha^{\mathrm{im}}}(w) \;:=\; \oint_{\gamma_w}\,V_{\alpha^{\mathrm{im}}}(z)\,\mathrm{d}z
  \;=\; \oint_{\gamma_w}\,{:}e^{\alpha^{\mathrm{im}}\cdot\phi(z)}{:}\,\mathrm{d}z,
\end{equation}
with $\gamma_w$ a small loop around $w$ chosen so that
$V_{\alpha^{\mathrm{im}}}(z)$ has conformal weight one. The OPE
$V_{\alpha^{\mathrm{im}}}(z)V_{\beta^{\mathrm{im}}}(w)$ has pole order
$-\langle\alpha^{\mathrm{im}},\beta^{\mathrm{im}}\rangle$; when
$\langle\alpha,\beta\rangle \geq 0$ the OPE is regular and the screening
commutator
$[S_{\alpha^{\mathrm{im}}}, S_{\beta^{\mathrm{im}}}] = 0$ in agreement
with Borcherds axiom (GKM4) (the imaginary-root vanishing axiom).
When $\langle\alpha,\beta\rangle = -1$, the bracket extracts a single
residue giving $e^{\mathrm{im}}_{\alpha+\beta}$ with coefficient
$c_{K3}(4(nm' + n'm) - 2\ell\ell')$. The structure constants
$c_{\alpha,\beta}^j$ of the imaginary-root bracket
\begin{equation}
\label{eq:voa-imag-bracket}
  [e^{\mathrm{im}}_\alpha, e^{\mathrm{im}}_\beta] \;=\; \begin{cases}
    0, & \langle\alpha,\beta\rangle \geq 0, \\
    \sum_j c_{\alpha,\beta}^{\;j}\,e^{\mathrm{im}}_{\alpha+\beta - \gamma_j}, & \langle\alpha,\beta\rangle < 0,
  \end{cases}
\end{equation}
are determined by the K3 elliptic-genus Fourier coefficients at the
pair weight, via Borcherds' singular theta lift (Borcherds 1998,
\S6).

\paragraph{Weight-one primary condition and ghost cancellation.}
For $V_{\alpha^{\mathrm{im}}}$ to contribute a weight-one primary
(so that the contour integral is well-defined and BRST-closed), the
conformal weight
$h(V_{\alpha^{\mathrm{im}}}) = \tfrac{1}{2}\langle\alpha^{\mathrm{im}},\alpha^{\mathrm{im}}\rangle$
must be corrected by a transverse $c = 23$ SCFT sector and a single
$bc$-ghost pair at conformal weight $(\tfrac{1}{2},\tfrac{1}{2})$,
producing total ghost central charge $c_{\mathrm{gh}} = -26$. The
ghost-dressed imaginary-root vertex operator
\begin{equation}
\label{eq:voa-ghost-dressed}
  \tilde V_{\alpha^{\mathrm{im}}}(z) \;=\; V_{\alpha^{\mathrm{im}}}(z) \;\otimes\; c(z),
\end{equation}
after integration against $\mathrm{d}z$, lands in BRST cohomology
$H^\bullet(Q_{\mathrm{BRST}})$ and descends to
$e^{\mathrm{im}}_\alpha \in \mathfrak g_{\Delta_5}$.

\paragraph{The Borcherds denominator $1/\Delta_5$.}
The Weyl--Kac--Borcherds denominator function of $\mathfrak g_{\Delta_5}$
is (Gritsenko--Nikulin 1998, Theorem 2.1)
\begin{equation}
\label{eq:voa-GN-denom}
  \Delta_5(Z) \;=\; \sum_{(n,\ell,m) > 0}\,c_{\phi_{-2,1}}(4nm - \ell^2)\,q^n r^\ell s^m,
  \qquad q = e^{2\pi i\tau},\;r = e^{2\pi i z},\;s = e^{2\pi i\tau'},\; Z = (\tau,z,\tau'),
\end{equation}
the Gritsenko cusp form of weight $5$ on the Siegel upper half-space
$\Hbb_2$ with character $\chi_5$. The character is a square root of
the Igusa cusp form: $\Phi_{10}(Z) = \Delta_5(Z)^2$ (Gritsenko--Nikulin 1998,
Eq.~4.6). The Weyl--Kac identity reads
\begin{equation}
\label{eq:voa-WKB-denom}
  \Delta_5(Z) \;=\; e^{2\pi i\langle \rho,\, Z\rangle}\,
  \prod_{\alpha \in \Delta_+(\mathfrak g_{\Delta_5})}
  \bigl(1 - e^{2\pi i\langle\alpha,\,Z\rangle}\bigr)^{\mathrm{mult}(\alpha)},
\end{equation}
with $\rho = \tfrac{1}{2}(\delta_1 + \delta_2 + \delta_3) \in
\Lambda_{\Delta_5}\otimes\Qbb$ the Weyl vector of $\mathfrak g_{\Delta_5}$
at pairing $\langle\rho,\delta_i\rangle = -1$ for each real simple
root (half-sum of positive real roots, Gritsenko--Nikulin 1998
\S2, equation 2.5), and $\mathrm{mult}(\alpha)$ given by
$c_{K3}(\tfrac{1}{2}\langle\alpha,\alpha\rangle)$ (real root:
$\mathrm{mult} = 1$; imaginary root:
$\mathrm{mult} = c_{K3}(D)$ at discriminant $D$). The lowest-order
imaginary-root Fourier coefficients populating \eqref{eq:voa-imag-mult}
are $c_{K3}(0) = 2$, $c_{K3}(1) = 90$, $c_{K3}(2) = 462$,
$c_{K3}(3) = 1540$, $c_{K3}(4) = 4554$ (K3 elliptic genus
Eichler--Zagier 1985 Table 9). The generating function $1/\Delta_5$
is thus the partition function of the universal
$\mathfrak g_{\Delta_5}$-module $\bfH_{\Delta_5}$ in the natural
grading:
\begin{equation}
\label{eq:voa-partition-H-Delta5}
  \mathrm{ch}\,\bfH_{\Delta_5}(Z) \;=\; \frac{1}{\Delta_5(Z)}
  \;=\; e^{-2\pi i\langle\rho,\,Z\rangle}\,
  \prod_{\alpha\in\Delta_+}\,\bigl(1 - e^{2\pi i\langle\alpha,\,Z\rangle}\bigr)^{-\mathrm{mult}(\alpha)}.
\end{equation}
This is the direct Fock-space counting of states of $\bfH_{\Delta_5}$,
weighted by their $\Lambda_{\Delta_5}$ charge.

\paragraph{Two $\kappa$-invariants on $\bfH_{\Delta_5}$: chiral versus BKM.}
The K3 chiral bialgebra $\bfH_{\Delta_5}$ carries two distinct
$\kappa$-invariants, which must not be conflated. The chiral
$\kappa_{\mathrm{ch}}(\bfH_{\Delta_5}) = 3$ equals the positive
rank of the K3 transcendental lattice minus the Mukai-doubling
defect:
$\kappa_{\mathrm{ch}} = \rho(K3) - 1 + \chi(\cO_{K3})/2 = 1 + 2 = 3$
for a polarised K3 at Picard rank one (the programme-canonical K3 of
Vol III Theorem~CY-A$_3$). The BKM-weight invariant
$\kappa_{\mathrm{BKM}}(\Phi_1) = c_1(0)/2 = 5$ equals the weight of
the Gritsenko--Nikulin form $\Delta_5$ at the universal
eight-form-spread index $N = 1$, with $c_1(0) = 10$ the Fourier
zero-coefficient. The two are distinct mathematical objects living
on different complexes: $\kappa_{\mathrm{ch}}$ is the chiral
homology shadow computed from the bar complex of $A_{K3\times E}$,
$\kappa_{\mathrm{BKM}}$ is the Borcherds singular-theta lift weight
computed from the K3 elliptic-genus Jacobi form
$\phi^{K3}_{-2,1}$. The additive split
$\kappa_{\mathrm{BKM}} = \kappa_{\mathrm{ch}} + \chi(\cO_{\mathrm{fiber}})$
is false: the fiber $\cO_E$ has $\chi = 0$, giving $3 + 0 = 3 \neq 5$.
The correct relation is the master universal Borcherds weight
identity $\kappa_{\mathrm{BKM}}(\Phi_N) = c_N(0)/2$ (Borcherds 1995
singular-theta lift, Gritsenko 1999 explicit denominator formula).

% ==================================================
\subsection{Universal trace $\hbar^2 K = -1$ on K3: Bruinier 2002 Proposition 5.1 and Mukai signature}
\label{subsec:voa-hbar-K-universal-trace}
% ==================================================

The master trace identity $\hbar^2 K = -1$ of
Theorem~\ref{thm:climax-IV} on $X = K3$ reduces to the
representation-theoretic identification $K = 2c_+(L) = 8$ via
Bruinier 2002 Proposition 5.1 and the Serre-duality normalisation
of the vacuum trace on the derived centre.

\paragraph{Bruinier Heegner--Chern class reciprocity (Prop.~5.1).}
For an even lattice $L$ of signature $(2, \ell)$ and a Heegner
divisor $H_\beta(m) \subset \cF_L = \mathrm{O}^+(L)\backslash\cH_L$
(with $\cH_L$ the period domain, $\beta \in L^\vee/L$ the coset, and
$m < 0$ the norm), Bruinier 2002 Proposition 5.1 gives the Chern-class
identity
\begin{equation}
\label{eq:voa-bruinier-heegner-chern}
  c_1\bigl(\cL_{\Psi}\bigr) \;=\; \frac{k}{2}\,c_1(\mathcal L)
  \;-\; \sum_{\substack{\beta \in L^\vee/L \\ m < 0}}\,c_\Psi(\beta, m)\,[H_\beta(m)],
\end{equation}
with $\cL$ the automorphic line bundle on $\cF_L$, $\Psi$ a Borcherds
product of weight $k = c_\Psi(0, 0)/2$, and $c_\Psi(\beta, m)$ the
principal-part Fourier coefficients of the input weakly-holomorphic
modular form $F_\Psi$. The quantity $c_+(L)$ is the positive-signature
rank of $L$, so $c_+(\mathrm{Muk}(K3)) = 4$ and $K = 2c_+(L) = 8$ in
every enhancement compatible with the Mukai lattice.

\paragraph{Mukai signature and the $K = 8$ identification.}
The Mukai lattice of K3, $L = H^\bullet_{\mathrm{even}}(K3, \Zbb) =
H^0 \oplus H^2 \oplus H^4$ with Mukai pairing
$\langle\alpha,\beta\rangle_{\mathrm{Muk}} = \int_{K3}\alpha\cup\iota(\beta)$
(with $\iota$ the Mukai involution), has signature $(4, 20)$: four
positive-definite classes (degree-zero and degree-four classes $H^0
\oplus H^4 = \Zbb^2$ together with $\mathrm{Pic}(K3) \cap
(H^{1,1}(K3,\Rbb))^+$ of rank 2 at the Mukai-lattice basepoint) and
twenty negative-definite classes (the transcendental lattice of K3
together with the remaining positive $(1,1)$-classes). The
Mukai-doubling identification of Vol III (Theorem~7.3 of
\texttt{k3\_chiral\_bialgebra\_platonic.tex} in Vol III) gives
$c_+(L) = 4$, hence
\begin{equation}
\label{eq:voa-K-8-Mukai}
  K \;=\; 2\,c_+(L) \;=\; 8, \qquad \hbar^2 \;=\; -1/K \;=\; -1/8, \qquad \hbar^2 K \;=\; -1.
\end{equation}

\paragraph{Serre-duality normalisation of the vacuum trace.}
The Calabi--Yau-$3$ structure on $K3 \times E$ (Kontsevich--Vlassopoulos
2022, Theorem 1.2) endows the derived centre $Z^{\mathrm{der}}_{\mathrm{ch}}(A_{K3\times E})$
with a trace pairing $\tr : Z^{\mathrm{der}}_{\mathrm{ch}} \to \Cbb[-3]$
satisfying $\tr \circ \mathbb S_{D^b(X)} = -\tr$: the Serre functor
inverts the trace. On the vacuum module $\Omega_A \in A\text{-mod}$,
Caldararu--Willerton 2010 Theorem 1.6 motivates the comparison
\begin{equation}
\label{eq:voa-trace-vacuum}
  \tr(1_A) \;=\; \hbar^2\cdot K,
\end{equation}
with the equality read as a normalized conductor identity. It is not
derived from the Todd integral of a trivial rank-$K$ object on
$K3\times E$: that Euler characteristic is zero. The
$\mathbb S$-anti-self-duality fixes the sign convention for the trace;
the value $\hbar^2K=-1$ is the imposed Mukai--Bruinier conductor
normalisation.

\paragraph{Anomaly off the Koszul locus.}
Off the Koszul locus, Gerstenhaber's cyclic-pairing obstruction
produces the anomalous identity
\begin{equation}
\label{eq:voa-universal-trace-anomaly}
  \hbar^2 K \;+\; \int_X\,\at_X\cup B^{(2)} \;=\; -1,
\end{equation}
with the correction term given by the cocycle of
Construction~\ref{constr:cocycle}. On $\cU^{\mathrm{adm}}(K3)$ the
integral vanishes and the identity reduces to $\hbar^2 K = -1$
exactly; on the generic quintic $Q_5$ the integral contributes an
irreducible $101$-dimensional shift and the master identity acquires
an anomaly incompatible with $K \in \Zbb$.

\paragraph{Explicit target of $\at_X \cup B^{(2)}$ on K3 and $K3\times E$.}
On $X = K3$ the vanishing of $\int_X\,\at_X\cup B^{(2)}$ is a
direct consequence of the Mukai-lattice Atiyah computation
(Kapranov 1999 Prop.~4.4, Remark): the tangent bundle $T_{K3}$
has $\at_{K3} \in \Ext^1_{K3}(T_{K3}, T_{K3}\otimes\Omega^1_{K3})$,
and on a K3 (Calabi--Yau surface) the Dolbeault representative is
$\bar\partial\nabla \in \Omega^{0,1}(K3, \End(T_{K3})\otimes\Omega^1_{K3})$
for a Chern connection. The cup product $\at_{K3}\cup B^{(2)}$
with Connes' bidegree-two cyclic shift takes values in
$H^2(K3, \Omega^1_{K3})=H^{1,2}(K3)$, hence vanishes by dimension.
This proves only the K3-only receptacle vanishing. For $K3\times E$,
the group $H^2(K3\times E,\Omega^1)$ is nonzero of dimension $21$;
vanishing there must be a computation of the specific projected class,
not vanishing of the receptacle.

% ==================================================
\subsection{Theorem B chiral Positselski: explicit bar-cobar on $\cH_1$ and $V_1(\mathfrak{sl}_2)$}
\label{subsec:voa-theorem-B-bar-cobar}
% ==================================================

The canonical non-trivial witnesses of Theorem B on the reference
curve $C$ are the rank-one Heisenberg $\cH_1$ (abelian OPE, simplest
possible) and the level-one affine $V_1(\mathfrak{sl}_2)$ (non-abelian
OPE, simplest possible with simple-pole structure).

\paragraph{The Heisenberg $\cH_1$: abelian OPE, trivial bar differential.}
Single chiral generator $\alpha(z)$ of conformal weight one, OPE
\begin{equation}
\label{eq:voa-heis-ope}
  \alpha(z)\,\alpha(w) \;=\; \frac{1}{(z-w)^2} \;+\; \mathrm{regular}, \qquad
  [\alpha_m, \alpha_n] \;=\; m\,\delta_{m+n, 0}\cdot\mathbf{1},
\end{equation}
central charge $c = 1$, $\kappa(\cH_1) = 1$. Abelian: double pole only,
no simple pole. The augmentation ideal
$\bar\cH_1 = \bigoplus_{n \geq 1}\Cbb\cdot{:}\partial^{k_1}\alpha\cdots\partial^{k_n}\alpha{:}$
generates the ordered bar complex
$\Bar^{\mathrm{ch}}_C(\cH_1) = T^c_C(s^{-1}\bar\cH_1)$ with
Maurer--Cartan differential driven by pair residues.

\emph{Bar differential on pairs.} The Borcherds identity
(Borcherds 1986; see $\S$\ref{sub:B2-operad}) gives
\begin{equation}
\label{eq:voa-heis-bar-diff}
  d_{\mathrm{res}}\bigl[s^{-1}\alpha(z)\,|\,s^{-1}\alpha(w)\bigr]
  \;=\; \mathrm{Res}_{z = w}\bigl((z-w)\cdot (z-w)^{-2}\bigr)\cdot\mathbf{1}
  \;=\; 0.
\end{equation}
The simple-pole coefficient $\alpha_{(0)}\alpha = 0$ is the Heisenberg
complementarity conductor $K_{\mathrm{Heis}} = 0$. Higher bar
differentials concatenate residues of the vanishing simple pole and
vanish identically.

\emph{Cobar differential and Positselski contracting homotopy.}
$\Cobar^{\mathrm{ch}}_C(\Bar^{\mathrm{ch}}_C(\cH_1)) = T(s\overline{\Bar^{\mathrm{ch}}(\cH_1)})$
carries the cobar cocomposition plus the (trivial) inherited bar
differential. Positselski's chain homotopy
\begin{equation}
\label{eq:voa-heis-homotopy}
  h\,:\, \Cobar\,\Bar \;\longrightarrow\; \Cobar\,\Bar[-1], \qquad
  h\bigl(s^{-1}x \otimes s\,y\bigr) \;=\; x\cdot y
\end{equation}
is the tautological shift (no curvature correction because the bar
differential is null). One computes
$[d, h] = d\circ h + h\circ d = \id - p_0$, with $p_0$ the projection
onto the weight-zero (vacuum) line. Positselski's Theorem 6.5 gives
the chain-level equivalence
\begin{equation}
\label{eq:voa-heis-inversion}
  \epsilon : \Cobar^{\mathrm{ch}}_C\,\Bar^{\mathrm{ch}}_C(\cH_1) \;\xrightarrow{\sim}\; \cH_1, \qquad
  [x_1|\cdots|x_n] \;\longmapsto\; {:}x_1\cdots x_n{:},
\end{equation}
the counit being normal-ordered reconstruction. All obstruction
classes $\mathrm{obs}_n \in H^n(\cH_1, \cH_1 \otimes \Omega^1_C)$
vanish for $n \geq 2$; $\cH_1$ is Koszul on the nose.

\emph{Explicit weight-two counit on Fock states.} The weight-two
Fock states of $\cH_1$ are
$\alpha_{-1}^2\Omega$ and $\alpha_{-2}\Omega$ with inner products
$\langle\alpha_{-1}^2\Omega|\alpha_{-1}^2\Omega\rangle = 2$,
$\langle\alpha_{-2}\Omega|\alpha_{-2}\Omega\rangle = 2$,
$\langle\alpha_{-1}^2\Omega|\alpha_{-2}\Omega\rangle = 0$. The counit
\eqref{eq:voa-heis-inversion} at weight two reads
\[
  \epsilon\bigl([s^{-1}\alpha \otimes s^{-1}\alpha]\bigr) \;=\; {:}\alpha\alpha{:}(0)\Omega \;=\; \alpha_{-1}^2\Omega,
  \qquad \epsilon\bigl([s^{-1}\partial\alpha]\bigr) \;=\; \partial\alpha(0)\Omega \;=\; \alpha_{-2}\Omega,
\]
with no cross-contamination between the two modes (orthogonality of
$\alpha_{-1}^2$ and $\alpha_{-2}$ on the vacuum). The Connes cyclic
shift $B^{(2)}$ at bidegree two acts on the cyclic tensor
$[s^{-1}\alpha | s^{-1}\alpha]$ as the $\Zbb/2$-antisymmetrisation,
which vanishes by $\alpha\otimes\alpha = \alpha\otimes\alpha$
(abelian commutation). Hence $B^{(2)}$ annihilates every weight-two
cyclic bar chain on $\cH_1$, and the cocycle $[m_3, B^{(2)}]_{\cH_1}
= 0$ on the nose in all homological degrees.

\paragraph{The level-one affine $V_1(\mathfrak{sl}_2)$: Frenkel--Kac realisation.}
Currents $J^a(z)$ for $a \in \{+,-,0\}$, level $k = 1$, OPE
\begin{equation}
\label{eq:voa-V1sl2-ope}
  J^a(z)\,J^b(w) \;=\; \frac{\delta^{ab}}{(z-w)^2} \;+\; \frac{i\,\varepsilon^{abc}\,J^c(w)}{z - w}
  \;+\; \mathrm{regular},
\end{equation}
central charge $c = \dim(\mathfrak{sl}_2)\cdot k/(k + h^\vee) =
3\cdot 1/(1+2) = 1$, $\kappa(V_1(\mathfrak{sl}_2)) =
\dim(\mathfrak{sl}_2)(k + h^\vee)/(2 h^\vee) = 3(1+2)/4 = 9/4$.

\emph{Frenkel--Kac boson realisation.} Frenkel--Kac 1980 identifies
$V_1(\mathfrak{sl}_2) \cong V_{A_1}$, the rank-one lattice vertex
algebra on the root lattice $A_1 = \Zbb\alpha$ with $\langle\alpha,\alpha\rangle = 2$,
through the map
\begin{equation}
\label{eq:voa-FK-realisation}
  J^0(z) \;=\; \frac{1}{\sqrt{2}}\,\partial\phi(z), \qquad
  J^{\pm}(z) \;=\; {:}\exp\bigl(\pm\sqrt{2}\,\phi(z)\bigr){:} \cdot \epsilon^\pm,
\end{equation}
with $\phi$ a free boson, $[\phi_m, \phi_n] = m\delta_{m+n,0}$, and
$\epsilon^\pm$ cocycle factors resolving the $\mathbb Z/2$-gradation
ambiguity. The lattice-vertex-operator OPE (Kac 1998, \S5.5) gives
\begin{equation}
\label{eq:voa-FK-ope-check}
  V_\alpha(z)\,V_{-\alpha}(w) \;=\; (z-w)^{-2}\cdot\mathbf{1} + (z-w)^{-1}\,\sqrt{2}\,\partial\phi(w) + \mathrm{reg.},
\end{equation}
and comparison with \eqref{eq:voa-V1sl2-ope} identifies
$V_\alpha = J^+$, $V_{-\alpha} = J^-$, $\sqrt{2}\,\partial\phi = 2\,J^0$.

\emph{Bar differential: Kac--Moody bracket on desuspensions.}
Pair-level bar differential
\begin{equation}
\label{eq:voa-V1sl2-bar-diff}
  d_{\mathrm{res}}\bigl[s^{-1}J^a\,|\,s^{-1}J^b\bigr]
  \;=\; \mathrm{Res}_{z = w}\Bigl(\frac{i\,\varepsilon^{abc}\,J^c(w)}{z - w}\Bigr)\,
  \mathrm{d}(z - w) \;=\; i\,\varepsilon^{abc}\cdot s^{-1}J^c,
\end{equation}
extracting the Kac--Moody bracket, plus double-pole Koszul
corrections from the $\delta^{ab}$ central term.

\emph{Bar cohomology.} By Positselski 2011 \S6.5 applied to
$\hat{\mathfrak{sl}}_2$ at level one, the bar cohomology computes
\begin{equation}
\label{eq:voa-V1sl2-bar-cohomology}
  H^\bullet\,\Bar^{\mathrm{ch}}_C(V_1(\mathfrak{sl}_2))
  \;=\; \mathrm{Sym}^\bullet(s^{-1}\mathfrak{sl}_2)\otimes\Lambda^\bullet(s^{-1}\Omega^1_C),
\end{equation}
the Koszul resolution of the universal enveloping chiral algebra.
The obstruction tower collapses at $n = 2$ via the Jacobi identity:
$\mathrm{obs}_2(J^a, J^b, J^c) = J^a[J^b, J^c] + \text{cyc.} = 0$.
$V_1(\mathfrak{sl}_2)$ sits in the Koszul stratum.

\emph{Counit and Wakimoto resolution.} The Positselski counit
\begin{equation}
\label{eq:voa-V1sl2-counit}
  \epsilon : \Cobar^{\mathrm{ch}}_C(\Bar^{\mathrm{ch}}_C(V_1(\mathfrak{sl}_2))) \;\longrightarrow\; V_1(\mathfrak{sl}_2),
  \qquad [J^{a_1}|\cdots|J^{a_n}] \;\longmapsto\; {:}J^{a_1}(z)\cdots J^{a_n}(z){:}
\end{equation}
is a quasi-isomorphism in $D^{\mathrm{co}}_{\mathrm{ch}}(C)$; the
inverse-modulo-homotopy is witnessed by the Wakimoto resolution
(Frenkel--Ben-Zvi 2004, \S15.4): every $V_1(\mathfrak{sl}_2)$-module
$M$ admits a resolution
\begin{equation}
\label{eq:voa-wakimoto}
  0 \;\longrightarrow\; M \;\longrightarrow\; W_0 \;\xrightarrow{d_0}\; W_1 \;\xrightarrow{d_1}\; \cdots
\end{equation}
by semi-infinite Fock modules $W_j$ over the sub-Heisenberg
$\Cbb\cdot\partial\phi(z)$; each $W_j$ is Koszul on the nose
(as in the $\cH_1$ computation above); the hypercohomology spectral
sequence converges to $M$ with trivial $E_2$-page after the first
differential. The cocycle $[m_3, B^{(2)}]$ restricts to zero on
$V_1(\mathfrak{sl}_2)$ because $\End^\bullet(V_1(\mathfrak{sl}_2))$ is
strictly associative: $m_3 \equiv 0$ as an operadic chain,
Massey-bracket obstructions vanish per level-one-triviality of the
$A_\infty$-lift.

\paragraph{Joint conclusion.}
Both $\cH_1$ and $V_1(\mathfrak{sl}_2)$ witness Theorem B chain-level
with explicit contracting homotopy; $[m_3, B^{(2)}]$ vanishes on the
reference-curve shadow of each; the comparison targets specialise to
identities on Fock and lattice-vertex states.

% ==================================================
\subsection{3D gravity $1/\Phi_{10}$: DVV 1997 derivation via symmetric-orbifold VOA on $\mathrm{Sym}^N(K3)$}
\label{subsec:voa-dvv-symorbifold}
% ==================================================

The Dijkgraaf--Moore--Verlinde--Verlinde 1997 derivation of
$1/\Phi_{10}$ as a generating function of K3 BPS microstates proceeds
through the symmetric-orbifold vertex operator algebra.

\paragraph{Symmetric-orbifold VOA on $\mathrm{Sym}^N(K3)$.}
For the K3 sigma-model superconformal field theory with
$(c,\bar c) = (6,6)$ central charges, let $\cV_{K3}$ denote the
corresponding vertex operator super-algebra with elliptic-genus
character
\begin{equation}
\label{eq:voa-K3-ell-genus}
  \chi_{\mathrm{ell}}(K3; \tau, z) \;=\; 8\,\Bigl[
    \Bigl(\frac{\theta_2(\tau,z)}{\theta_2(\tau,0)}\Bigr)^2 +
    \Bigl(\frac{\theta_3(\tau,z)}{\theta_3(\tau,0)}\Bigr)^2 +
    \Bigl(\frac{\theta_4(\tau,z)}{\theta_4(\tau,0)}\Bigr)^2
  \Bigr] \;=\; 2\phi_{0,1}(\tau,z),
\end{equation}
the weak Jacobi form of weight $0$ and index $1$ (Eichler--Zagier 1985,
\S9). The $N$-fold symmetric orbifold $\mathrm{Sym}^N(K3) = K3^N/S_N$
carries the symmetric-orbifold VOA
\begin{equation}
\label{eq:voa-sym-orbifold}
  \cV_{\mathrm{Sym}^N(K3)} \;=\; \bigl(\cV_{K3}^{\otimes N}\bigr)^{S_N} \;\oplus\; \bigoplus_{[\sigma]}\cV_{K3}^{(\sigma)},
\end{equation}
with $[\sigma]$ running over non-identity conjugacy classes of $S_N$
and $\cV_{K3}^{(\sigma)}$ the twisted sector associated to each
cycle-type partition of $N$. The untwisted sector is the $S_N$-invariant
tensor power; the twisted sectors contribute one copy of $\cV_{K3}$
per cycle in the partition, compactified on the long cycle.

\emph{Twisted-sector generators at small $N$.} For $N = 2$, the
non-trivial conjugacy class is the $2$-cycle, and the twisted sector
$\cV_{K3}^{(12)}$ is a single copy of $\cV_{K3}$ compactified on a
length-$2$ cycle, i.e.\ states live on the double cover
$\tilde S^1 \to S^1$ with period $4\pi$ on the double cover versus
$2\pi$ on the base. The twisted-sector ground state is a rank-one
module over the $\Zbb/2$-invariant subalgebra
$\cV_{K3}^{S_2} = \cV_{K3}\otimes\cV_{K3}$ with twisted vertex operators
\[
  V_{(12), \alpha\otimes\beta}(z) \;=\; \frac{1}{\sqrt{2}}\bigl(V_\alpha(z)\cdot V_\beta(-z) + V_\beta(z)\cdot V_\alpha(-z)\bigr),
\]
for $\alpha,\beta \in L_{K3}$, the Lehn--Sorger 2003 twisted vertex
operators on the Hilbert scheme. For $N = 3$, the conjugacy classes are
$3$-cycles (one class) and pairs of transpositions (absent in $S_3$);
the $3$-cycle twisted sector is $\cV_{K3}$ on the triple cover, with
ground state partition function $q^{c/24}$ fraction at $c_{K3} = 6$
giving $q^{1/4}$ shift per the orbifold-GSO sum.

\paragraph{DMVV product formula.}
Dijkgraaf--Moore--Verlinde--Verlinde 1997 Theorem 5.1 computes the
symmetric-orbifold elliptic-genus generating function:
\begin{equation}
\label{eq:voa-dmvv}
  \sum_{N \geq 0}\,p^N\,\chi_{\mathrm{ell}}(\mathrm{Sym}^N(K3); \tau, z)
  \;=\; \prod_{n \geq 1, m \geq 0, \ell \in \Zbb}
  \frac{1}{(1 - p^n q^m r^\ell)^{c(nm, \ell)}},
\end{equation}
with $c(D, \ell)$ the Fourier coefficient of
$\chi_{\mathrm{ell}}(K3; \tau, z) = \sum_{m,\ell}c(m,\ell)q^m r^\ell$.
Under the multiplicative lift (Borcherds 1995; Gritsenko 1999), the
right-hand side equals the Gritsenko lift of $\phi^{K3}_{0,1}$:
\begin{equation}
\label{eq:voa-gritsenko-lift}
  \prod_{n,m,\ell}(1 - p^n q^m r^\ell)^{c(nm,\ell)}
  \;=\; p^{-1}q^{-1}r^{-1}\cdot\Phi_{10}(Z)
  \quad \text{up to the Gritsenko leading term,}
\end{equation}
with $\Phi_{10}$ the Igusa cusp form of weight $10$ on the Siegel upper
half-space (Igusa 1962; Gritsenko--Nikulin 1998, \S2). Combining
\eqref{eq:voa-dmvv} and \eqref{eq:voa-gritsenko-lift},
\begin{equation}
\label{eq:voa-dvv-1-over-Phi10}
  \sum_{N \geq 0}\,p^N\,\chi_{\mathrm{ell}}(\mathrm{Sym}^N(K3); \tau, z)
  \;=\; \frac{p q r}{\Phi_{10}(Z)}
  \;\propto\; \frac{1}{\Phi_{10}(Z)}.
\end{equation}

\paragraph{Identification with the 3D gravity partition function.}
Maldacena--Moore--Strominger 1999 identify the left-hand side of
\eqref{eq:voa-dvv-1-over-Phi10} with the $1/4$-BPS microstate-counting
partition function of type IIB string theory on $K3 \times S^1$ at
the attractor point. The Strominger--Vafa 1996 AdS$_3$/CFT$_2$
duality then identifies this CFT partition function with the
twisted-holography boundary character of 3D quantum gravity on
$\mathrm{AdS}_3 \times K3$:
\begin{equation}
\label{eq:voa-3dqg-partition}
  Z^{\AdS_3 \times K3}_{\mathrm{3dQG, BPS}}(Z) \;=\;
  \sum_{N \geq 0}\,p^N\,\chi_{\mathrm{ell}}(\mathrm{Sym}^N(K3); \tau, z)
  \;\propto\; \frac{1}{\Phi_{10}(Z)}.
\end{equation}
The boundary CFT is the $\mathrm{Sym}^N(K3)$ orbifold sigma-model
(Seiberg--Witten 1999; Dijkgraaf 1999 review), whose VOA lives on the
asymptotic $S^1 \times \Rbb$ boundary of $\mathrm{AdS}_3$; the
elliptic-genus insertion computes the BPS trace. The $3$D gravity
identification is at the twisted-holography scope: the dynamical
metric path integral is conjectural (Maloney--Witten 2009
sum convergence plus Maxfield--Turiaci 2020 wormhole corrections),
but the automorphic function $1/\Phi_{10}$ is rigorously pinned by
\eqref{eq:voa-dvv-1-over-Phi10} and the identification
$\Phi_{10} = \Delta_5^2$ of Gritsenko--Nikulin 1998.

% ==================================================
\subsection{Candidate arithmetic VOAs $V^{\mathrm{arith}}_{K3, d_K}$ at the conditional Heegner discriminants}
\label{subsec:voa-heegner-arithmetic}
% ==================================================

Each of the nine Heegner discriminants
$d_K \in \{3, 4, 7, 8, 11, 15, 19, 20, 24\}$ of
Proposition~\ref{prop:heegner} is treated here as a conditional
candidate for an arithmetic VOA
$V^{\mathrm{arith}}_{K3, d_K}$ with central charge $K = 8$. The
canonical construction is not proved in this supplement; it requires
the rank-$8$ Bruinier coefficient computation and an explicit
class-group action on the Mukai-lattice VOA.

\paragraph{Imaginary-quadratic orders and Heegner CM points.}
For each $d_K$ above, let $K_{d_K} = \Qbb(\sqrt{-d_K})$ be the
imaginary-quadratic field of discriminant $-d_K$ and
$\cO_{d_K}$ the corresponding quadratic order. The values
$d_K=8,20,24$ are fundamental negative discriminants in the usual
field-discriminant convention here; they are not used as
conductor-$2$ replacements for $d_K=4,5,6$. The Heegner point
$\tau_{d_K} \in
\mathfrak H$ on the upper half-plane is the CM point of discriminant
$-d_K$ where the elliptic curve $E_{\tau_{d_K}} = \Cbb/\langle 1,
\tau_{d_K}\rangle$ has complex multiplication by $\cO_{d_K}$. Via
the Shioda--Inose 1977 correspondence, each ideal class of
discriminant $-d_K$ determines a singular K3 surface $K3_{d_K}$ with
transcendental lattice
$T_{d_K}$ of rank 2, positive-definite, discriminant $-d_K$:
\begin{equation}
\label{eq:voa-T-dK}
  T_{d_K} \;\cong\; \begin{cases}
    \Zbb\bigl\langle 1, \tfrac{1 + \sqrt{-d_K}}{2}\bigr\rangle, & d_K \equiv 3 \pmod 4, \\[4pt]
    \Zbb\bigl\langle 1, \tfrac{\sqrt{-d_K}}{2}\bigr\rangle, & d_K \equiv 0 \pmod 4,
  \end{cases}
\end{equation}
with Gram matrix
$G_{T_{d_K}} = \mathrm{diag}(2, (1 + d_K)/2)$ at odd $d_K$
and $\mathrm{diag}(2, d_K/2)$ at even $d_K$
(Shioda--Inose 1977, Proposition 4.3).

\paragraph{Scheithauer generalised-moonshine lift.}
At each $\tau_{d_K}$, the expected construction is a twisted VOA
$V^{\mathrm{arith}}_{K3, d_K}$ obtained from the Mukai-lattice VOA
$V_{\mathrm{Muk}(K3)}$ by a Scheithauer-type
generalised-moonshine lift and the class-group action of the order
$\cO_{d_K}$:
\begin{equation}
\label{eq:voa-Vari-K3-dK}
  V^{\mathrm{arith}}_{K3, d_K} \;:=\;
  \bigl(V_{\mathrm{Muk}(K3)}\bigr)^{\sigma_{d_K}}
  \otimes \Cbb[\mathrm{Cl}(\cO_{d_K})],
\end{equation}
where $\mathrm{Cl}(\cO_{d_K})$ is the ideal class group. The quotient
$\cO_{d_K}/\cO_{d_K}^*$ is not the class group and is not used.

\paragraph{Explicit generators at the class-number-one discriminants in the conditional list.}
For $d_K \in \{3, 4, 7, 8, 11, 19\}$ with class number $h_{d_K} = 1$,
the class group $\mathrm{Cl}(\cO_{d_K})$ is trivial and
$V^{\mathrm{arith}}_{K3, d_K} = (V_{\mathrm{Muk}(K3)})^{\sigma_{d_K}}$
has no non-trivial class-group twist. At each such $d_K$, the transcendental-lattice
generator $\theta_{d_K}(z) = V_{(1, \tau_{d_K})}(z)$ realises the
CM data:
\begin{itemize}
\item $d_K = 3$: $\theta_3(z) = V_{(1, \omega_3)}(z) = {:}e^{\omega_3\cdot\phi(z)}{:}$ with
$\omega_3 = (1 + \sqrt{-3})/2$; $\mathcal O_3 = \Zbb[\omega_3]$ Eisenstein integers.
\item $d_K = 4$: $\theta_4(z) = V_{(1, i)}(z)$ with $i = \sqrt{-1}$;
$\mathcal O_4 = \Zbb[i]$ Gaussian integers.
\item $d_K = 7$: $\theta_7(z) = V_{(1, \omega_7)}(z)$, $\omega_7 = (1 + \sqrt{-7})/2$;
$\mathcal O_7 = \Zbb[\omega_7]$.
\item $d_K = 8$: $\theta_8(z) = V_{(1, \sqrt{-2})}(z)$; $\mathcal O_8 = \Zbb[\sqrt{-2}]$.
\item $d_K = 11$: $\theta_{11}(z) = V_{(1, \omega_{11})}(z)$, $\omega_{11} = (1 + \sqrt{-11})/2$;
$\mathcal O_{11} = \Zbb[\omega_{11}]$.
\item $d_K = 19$: $\theta_{19}(z) = V_{(1, \omega_{19})}(z)$, $\omega_{19} = (1 + \sqrt{-19})/2$;
$\mathcal O_{19} = \Zbb[\omega_{19}]$.
\end{itemize}
Each $\theta_{d_K}$ has conformal weight
$h(\theta_{d_K}) = \tfrac{1}{2}\langle (1,\tau_{d_K}), (1,\tau_{d_K})\rangle$
computed from \eqref{eq:voa-T-dK}; the Scheithauer twist then
tensors with a rank-$8$ lattice-vertex-algebra shift landing
$V^{\mathrm{arith}}_{K3, d_K}$ at central charge $K = 8$.

The central-charge check at the fixed points: the $\sigma_{d_K}$-invariant
subalgebra of $V_{\mathrm{Muk}(K3)}$ has central charge equal to the
signature-positive rank $c_+(L) = 4$ (Mukai-doubling Vol III
\texttt{k3\_chiral\_bialgebra\_platonic.tex}); the Scheithauer
arithmetic twist doubles the central charge to $K = 2c_+(L) = 8$,
and the ghost dressing at level-one Heisenberg shifts by
$c_{\mathrm{gh}} = -26$ then back by $+26$ ($bc$ and $\beta\gamma$
complementary cancellation), netting the $K = 8$ invariant across
the nine strata.

\paragraph{Class-number-two discriminants in the conditional list.}
For $d_K \in \{15, 20, 24\}$ with $h_{d_K} = 2$, the class group
$\mathrm{Cl}(\cO_{d_K}) \cong \Zbb/2\Zbb$ is generated by a single
non-principal ideal class, and
$V^{\mathrm{arith}}_{K3, d_K} = (V_{\mathrm{Muk}(K3)})^{\sigma_{d_K}}
\otimes \Cbb[\Zbb/2]$ carries a candidate $\Zbb/2$-orbifold structure.
The assertion that exactly these class-number-two discriminants survive
a rank-$8$ Bruinier positivity cut is conditional; it is not a
consequence of Bruinier 2002 Proposition 5.1 alone. The other
class-number-two discriminants $d_K \in \{23, 31, 35, 39, 40, 43, 47, 51, 52, 55,
56, 68, 72, 84, 88, 100, 115, 120, 148, 168, 187, 195, 228, 232,
280, 292, 312, 340, 372, 403, 408, 427, 520, 532, 555, 595, 627,
708, 715, 760, 795, 1012, 1435\}$ produce Heegner divisors of wrong
signature for the rank-$8$ Mukai enhancement only after the missing
coefficient/sign calculation. At each admissible
$d_K$, $V^{\mathrm{arith}}_{K3, d_K}$ carries generators:
\begin{itemize}
\item $d_K = 15$: $\theta_{15}, \theta_{15}^\sigma$ under the non-principal
ideal $\mathfrak p_3 \cdot \mathfrak p_5$; $\mathcal O_{15} = \Zbb[(1+\sqrt{-15})/2]$.
\item $d_K = 20$: $\theta_{20}, \theta_{20}^\sigma$;
$\mathcal O_{20} = \Zbb[\sqrt{-5}]$.
\item $d_K = 24$: $\theta_{24}, \theta_{24}^\sigma$;
$\mathcal O_{24} = \Zbb[\sqrt{-6}]$.
\end{itemize}
The Borcherds--Kac--Moody algebra attached to
$V^{\mathrm{arith}}_{K3, d_K}$ is
$\mathfrak g^{\mathrm{arith}}_{d_K} := \mathfrak g_{\Delta_5}$ twisted
by $\sigma_{d_K}$ (Gritsenko--Nikulin 1998 \S7). Each
$V^{\mathrm{arith}}_{K3, d_K}$ has central charge $K = 8$ (Mukai
signature); the arithmetic twist alters the conformal weight
assignments of lattice vertex operators but leaves the central charge
and the $\kappa = 0$ Heisenberg conductor invariant.

\paragraph{Connection to the $\sfB$-row.}
Conditional on the arithmetic calculation in Proposition~\ref{prop:heegner} and on constructing the class-group action above, the nine $V^{\mathrm{arith}}_{K3, d_K}$ stratify the $\sfB$-row of the
derived-centre classification by arithmetic twist: each witnesses
the ceiling $K^\kappa = 8$ of Theorem~C with a different
imaginary-quadratic-field arithmetic. The CM point $\tau_{d_K}$ of
Heegner discriminant $-d_K$ sits on the Humbert divisor of
discriminant $d_K$ in $\cF_L$; the vanishing of
$[m_3, B^{(2)}]$ at $d_K$ is a scalar cyclic-HKR condition on the
corresponding Humbert-Heegner CM stratum, after embedding into the
principal $K3\times E$ CY3 specialization. Off the conditional list,
the expected obstruction is the missing rank-$8$ arithmetic
admissibility condition.

% ==================================================
\subsection{Koszul-dual algebras on the reference curve for the five vertices}
\label{subsec:voa-koszul-duals}
% ==================================================

Each of the comparison vertices of Theorem~\ref{thm:main} is expected to carry an
explicit Koszul-dual algebra $A^!$ on the reference curve $C$,
computed from $A^i = H^\star\Bar^{\mathrm{ch}}_C(A)$ via Verdier
duality $A^! = (A^i)^\vee$. The five Koszul duals assemble only after
the bridge hypotheses of Theorem~\ref{thm:main} are supplied:

\begin{itemize}
\item \emph{Vertex~$V_1 = \{[m_3, B^{(2)}]_X = 0\}$.} Koszul-dual of the
vanishing cocycle is the graded-commutative algebra
$\mathrm{Sym}^\bullet(H^{1,2}(X)^\vee[-1])$ on the dual of the
corrected obstruction target; trivial when $[m_3, B^{(2)}] = 0$ forces
$H^{1,2}(X)^\vee$ to act nilpotently on the chiral algebra
$A_X = \Phi_3(X)$ (via the Atiyah-class pairing).

\item \emph{Vertex~$V_2 = \mathrm{Thm\,B}$.} Koszul-dual of the chiral
Positselski inversion is the universal enveloping chiral coalgebra
$\Bar^{\mathrm{ch}}_C(A)$ itself on the cofree side; on the admissible
Koszul locus $A^! = \Omega^{\mathrm{ch}}_C(A^i)$ is the inversion
functor output. For $A = \cH_1$, $A^! = \Omega_C^\bullet$ the de Rham
complex on $C$ (trivial coalgebra on one generator). For
$A = V_1(\mathfrak{sl}_2)$, $A^! = \mathrm{CE}^\bullet(\hat{\mathfrak{sl}}_2)$
the Chevalley--Eilenberg complex computing the Lie-algebra cohomology.

\item \emph{Vertex~$V_3 = \mathrm{Stage\text{-}1\ canonicality}$.}
Koszul-dual of the canonicality of the stage-one
$E_3$-factorisation algebra $\Phi^{\mathrm{FA}}_3(X)$ is the stage-two
specialisation $\Sp^{\mathrm{ch}}_{\Sigma_{d-1}, C}$ as a
Koszul-functor; its Koszul dual on $C$ is the universal curve-side
deformation complex, which for $X = K3$ is the Hilbert-scheme chiral
algebra $\mathrm{Hilb}^n(K3)$-chiral-algebra shadow on $C$.

\item \emph{Vertex~$V_4 = \bfH_{\Delta_5}$.} Koszul-dual of
$\bfH_{\Delta_5}$ is the Feigin--Frenkel dual
$\bfH_{\Delta_5}^! = V_{\mathrm{Muk}(K3)}^\vee \otimes \Lambda^\bullet(\mathfrak g^{\mathrm{im}}_{\Delta_5})$,
the Koszul-dual lattice vertex algebra tensored with the exterior
algebra on the imaginary simple roots. Central charge of the dual:
$K^! = -K + 2\dim(\mathfrak h) = -8 + 2\cdot 3 = -2$; the
complementarity $K + K^! = 8 - 2 = 6$ is the Vol III $\sfB$-row
six-form-spread witness (off the principal $K + K^! = 8$ ceiling by
the rank-three Cartan shift).

\item \emph{Vertex~$V_5 = \{\hbar^2 K = -1\}$.} Koszul-dual of the
universal trace identity is the Hochschild cochain complex
$\HH^\bullet(A_X, A_X)$ computing the derived centre
$Z^{\mathrm{der}}_{\mathrm{ch}}(A_X)$; its Koszul dual on $C$ is the
curve-side evaluation of the Serre trace pairing
$\tr : Z^{\mathrm{der}}_{\mathrm{ch}}(A_X)|_C \to \Cbb[-3]$, landing
in $\Cbb[-3]$ for CY-3 and pinned at $\tr(1) = -1$ by Serre duality.
\end{itemize}

The resulting Verdier-dual decagon is a frontier conjecture: under
\textup{(B1)--(B4)} and the additional construction of these Koszul
duals, it should refine Conjecture~\ref{conj:atiyah-connes-pentagon}.

% =====================================================================
% End of VOA supplement
% =====================================================================
